\documentclass[11pt,a4paper]{article}

% ---------------------------------------------------------------------------
% Packages
% ---------------------------------------------------------------------------
\usepackage[margin=1in]{geometry}
\usepackage{amsmath,amssymb,amsthm}
\usepackage{graphicx}
\usepackage{booktabs}
\usepackage{hyperref}
\usepackage{xcolor}
\usepackage{algorithm}
\usepackage{algpseudocode}
\usepackage{natbib}
\usepackage{subcaption}
\usepackage{multirow}
\usepackage{url}

\hypersetup{
  colorlinks=true,
  linkcolor=blue,
  citecolor=blue,
  urlcolor=blue,
}

% ---------------------------------------------------------------------------
% Title block
% ---------------------------------------------------------------------------
\title{%
  \textbf{Latent Reservoir Scratchpads for LLMs:}\\
  \large Bidirectional Echo State Networks as External Working Memory\\
  and Cheap Latent Reasoning Substrates
}

\author{%
  Anonymous Authors
}

\date{March 2026}

% ---------------------------------------------------------------------------
% Macros
% ---------------------------------------------------------------------------
\newcommand{\LRS}{\textsc{LRS}}
\newcommand{\ESN}{\textsc{ESN}}
\newcommand{\RIL}{\textsc{RIL}}

\theoremstyle{definition}
\newtheorem{definition}{Definition}

% ===========================================================================
\begin{document}
\maketitle

\begin{abstract}
Large Language Models (LLMs) suffer from two fundamental limitations in
long-horizon tasks: (i) the KV-cache grows linearly with context length, and
(ii) multi-step symbolic reasoning requires recirculating hidden states through
full forward passes.  We propose the \textbf{Latent Reservoir Scratchpad
(\LRS)}: a fixed-size Echo State Network (ESN) integrated into a Transformer
as an external, bidirectional working memory.  The reservoir state $r_t \in
\mathbb{R}^{D_r}$ evolves at each token step via a cheap sparse
matrix-vector multiply; only the read/write \emph{interface} parameters are
trained.  We evaluate \LRS\ across three integration depths---bolt-on sidecar
(Track~A), surgically inserted layers (Track~B), and from-scratch hybrid
architecture (Track~C, RW-Transformer)---against compute-matched baselines
including YaRN context extension, Mamba-2, Infini-attention, and Titans.

\noindent\textbf{Negative-result statement.} In accordance with our
pre-registered proposal, we report all gate outcomes here regardless of
direction.  Gate~A (bolt-on sidecar) results are pending GPU evaluation; this
paper documents the codebase, methodology, and preliminary CPU-scale
experiments.  We commit to updating the paper with empirical results upon
completion of GPU training runs.
\end{abstract}

% ---------------------------------------------------------------------------
\section{Introduction}
\label{sec:intro}
% ---------------------------------------------------------------------------

Standard Transformer architectures \citep{vaswani2017attention} encode all
context in the KV-cache, which grows as $O(L \cdot d)$ with sequence length
$L$.  While efficient attention variants
\citep{dao2022flashattention,munkhdalai2024infini} partially address this
cost, they remain fundamentally bounded by the sequence length.  Recurrent
models such as Mamba \citep{gu2023mamba} maintain an $O(1)$ state but lack
precise content-addressable retrieval.

A complementary line of work, \emph{latent reasoning}, proposes recirculating
hidden states through the LLM's own forward pass as a reasoning substrate
\citep{hao2024coconut,geiping2025huginn}.  This gives additional ``thinking
time'' but at full-forward-pass cost per reasoning step and with potential
distributional drift since the LLM's embedding space was not designed for
recirculation.

\textbf{Our contribution.} We introduce the Latent Reservoir Scratchpad
(\LRS), which decouples the working memory substrate from the LLM.  A
fixed-weight Echo State Network \citep{jaeger2001echo} serves as:
\begin{enumerate}
  \item A \emph{fixed-size} continuous-state memory buffer with $O(1)$ memory
    footprint regardless of sequence length.
  \item A \emph{cheap latent reasoning substrate}: the reservoir can evolve
    for $K$ sub-steps between tokens via single sparse matrix-vector
    multiplies, providing thinking time at $\ll 1\%$ the cost of LLM forward
    passes.
\end{enumerate}
Only the read/write interface parameters (projections and optional LoRA
adapters) are trained; the reservoir weights $W_{\mathrm{res}}$ remain frozen,
providing bounded-state compression with guaranteed echo state property.

We further explore the \textbf{DeltaNet synergy hypothesis}: Qwen3.5's Gated
DeltaNet layers \citep{yang2024gated} use a delta-rule hidden state update
analogous to ESN recurrence, potentially reducing distributional mismatch when
the LLM reads reservoir states.

\paragraph{Negative-result commitment.}  As pre-registered in our
collaborative research proposal, we commit to publishing detailed findings
regardless of gate outcomes.  If any success gate fails, Section~\ref{sec:results}
documents the failure mode and analysis.

% ---------------------------------------------------------------------------
\section{Related Work}
\label{sec:related}
% ---------------------------------------------------------------------------

\subsection{Reservoir Computing}

Echo State Networks \citep{jaeger2001echo,jaeger2004harnessing} are recurrent
networks where only the output (readout) layer is trained; the reservoir
weights are fixed random sparse matrices with spectral radius near 1.
\citet{shen2021reservoir} inserted reservoir states into Transformers as
``Reservoir Transformers,'' finding modest gains on language modelling.
\citet{koester2025reservoir} demonstrated reservoir computing as a standalone
language model at small scale. The echo state property---that the reservoir
state depends only on recent inputs, not on the initial state---provides a
mathematical foundation for bounded memory.

\subsection{Memory Augmentation}

Infini-attention \citep{munkhdalai2024infini} and Titans
\citep{behrouz2025titans} augment Transformer layers with compressive
associative memory matrices.  Infini-attention maintains a per-layer
associative matrix updated with each segment; Titans introduces meta-learned
updates at test time.  Both require $O(d^2)$ memory per layer and $O(L \cdot
d^2)$ updates.  In contrast, \LRS\ maintains a \emph{single} reservoir state
of dimension $D_r$ with $O(D_r)$ memory and $O(D_r^2 / \rho_s)$ per-step
compute where $\rho_s \ll 1$ is reservoir sparsity.

\subsection{Latent Reasoning}

Coconut \citep{hao2024coconut} trains LLMs to reason by recirculating
continuous hidden states rather than discrete tokens.  Huginn
\citep{geiping2025huginn} loops a Transformer block for depth-adaptive
computation.  Ouro \citep{srivastava2025ouro} uses looped Transformers as
latent reasoning machines.  All three pay full forward-pass cost per
reasoning step.  Our reservoir sub-steps cost $O(D_r^2 \cdot \rho_s)$ per
step---orders of magnitude cheaper---while operating in a space $D_r \gg
d_{\mathrm{LLM}}$ for rich dynamics.

\subsection{Long-Context Transformers}

YaRN \citep{peng2023yarn} extends RoPE positional encodings to longer
contexts.  Mamba-2 \citep{gu2023mamba,dao2024transformers} replaces attention
with structured state-space layers for $O(L)$ sequence processing.  Our
evaluation includes both as baselines (Track A, Controls 1 and 2).

% ---------------------------------------------------------------------------
\section{Method}
\label{sec:method}
% ---------------------------------------------------------------------------

\subsection{Echo State Network Reservoir}
\label{sec:esn}

A reservoir is a high-dimensional dynamical system defined by:
\begin{equation}
  r_t = (1 - \alpha)\, r_{t-1}
        + \alpha\, \tanh\!\bigl(W_{\mathrm{res}}\, r_{t-1}
                                + W_{\mathrm{in}}\, x_t
                                + w_t\bigr)
  \label{eq:esn}
\end{equation}
where $r_t \in \mathbb{R}^{D_r}$ is the reservoir state, $\alpha \in (0,1]$
is the leak rate, $W_{\mathrm{res}} \in \mathbb{R}^{D_r \times D_r}$ is a
sparse random matrix with spectral radius $\rho < 1$ (Erdős–Rényi or
Watts-Strogatz topology), $W_{\mathrm{in}} \in \mathbb{R}^{D_r \times d}$ is
a fixed random input projection, and $w_t$ is the LLM's write signal.  All of
$W_{\mathrm{res}}$, $W_{\mathrm{in}}$ are fixed after initialisation---they
are never trained.

\subsection{Read/Write Interface}
\label{sec:interface}

The trainable interface consists of:
\begin{itemize}
  \item \textbf{WriteHead}: $w_t = \mathrm{Linear}(h_t) \in \mathbb{R}^{D_r}$,
    where $h_t$ is the LLM hidden state at the chosen insertion layer.
  \item \textbf{ReadProjection}: $m_t = \mathrm{Linear}(r_t) \in
    \mathbb{R}^{d_{\mathrm{LLM}}}$, mapping reservoir state back to LLM
    dimension.
  \item \textbf{CrossAttentionSidecar} (optional): LLM hidden $h$ as queries,
    reservoir states as keys/values.  Multi-head, with $Q$-norm for stability.
  \item \textbf{FiLM Modulation} (optional): gated feature-wise linear
    modulation of hidden states by reservoir state.
\end{itemize}

Gradient stops at the numpy/torch boundary: reservoir weights receive no
gradient signal.  Only interface parameters (and optional LoRA adapters) are
updated during training.

\subsection{Track A: Bolt-On Reservoir Sidecar}
\label{sec:track_a}

Track~A attaches a single ESN to a frozen Qwen3.5-0.8B-Base model.  Training
progresses in two phases:
\begin{enumerate}
  \item \textbf{Read-only}: Only ReadProjection is trained; WriteHead is
    bypassed ($w_t = 0$).  The reservoir evolves from $x_t = h_t$
    (last-layer hidden states) as input.  This validates that the LLM can
    extract useful information from reservoir states before learning to write
    to them.
  \item \textbf{Read/Write}: Both ReadProjection and WriteHead are trained.
    LoRA adapters (rank 16) are added to $\{q, k, v, o, \text{gate},
    \text{up}, \text{down}\}$ projections.
\end{enumerate}
Integration uses CrossAttentionSidecar, with reservoir states as K/V memory.
A hyperparameter sweep (Track A, T16) searches over reservoir size
$D_r \in \{500, 2000, 10000, 50000\}$, leak rate $\alpha \in \{0.1, 0.3,
0.7, 1.0\}$, spectral radius $\rho \in \{0.5, 0.9, 0.99, 1.1\}$, and
topology $\in \{\text{Erdős–Rényi}, \text{Watts-Strogatz}\}$.

\subsection{Track B: Reservoir Interaction Layers (RIL)}
\label{sec:track_b}

Track~B surgically inserts the \LRS\ into mid-to-late Transformer layers
(replacing selected Gated DeltaNet blocks in Qwen3.5).  A dual-timescale
multi-reservoir is introduced:
\begin{itemize}
  \item \textbf{Fast reservoir}: $\alpha_{\mathrm{fast}} \approx 0.9$,
    $\rho_{\mathrm{fast}} \approx 0.9$---short memory, responsive to recent
    tokens.
  \item \textbf{Slow reservoir}: $\alpha_{\mathrm{slow}} \approx 0.1$,
    $\rho_{\mathrm{slow}} \approx 0.5$---long memory, integrates context.
\end{itemize}
The combined read state $[r_{\mathrm{fast}}; r_{\mathrm{slow}}]$ is passed to
the cross-attention sidecar.  The DeltaNet synergy hypothesis predicts lower
perplexity degradation for Qwen3.5 (with DeltaNet layers) than for LLaMA-3.2
(pure softmax attention) when both are augmented identically.

\subsection{Track C: RW-Transformer (From-Scratch Hybrid)}
\label{sec:track_c}

Track~C trains a new $\sim$0.8B parameter model from scratch.  Each decoder
block contains three parallel branches fused by a learned gating layer:
\begin{enumerate}
  \item Local/global attention (exact pattern matching).
  \item Standard MLP (static knowledge).
  \item Dual multi-reservoir workspace (fast + slow ESN).
\end{enumerate}
The write path is optimised via straight-through estimators, and $K$-step
latent reasoning is evaluated with $K \in \{1, 2, 4, 8, 16\}$ sub-steps
between tokens.

\subsection{Multi-Step Latent Reasoning}
\label{sec:latent_reasoning}

Between token generation steps, the reservoir evolves autonomously:
\begin{equation}
  r_t^{(k)} = (1 - \alpha)\, r_t^{(k-1)}
             + \alpha\, \tanh\!\bigl(W_{\mathrm{res}}\, r_t^{(k-1)}\bigr),
  \quad k = 1, \ldots, K
  \label{eq:substep}
\end{equation}
with $r_t^{(0)}$ set after the write injection.  Each sub-step costs one
sparse matrix-vector multiply---$O(\rho_s D_r^2)$ FLOPs vs.\
$O(d_{\mathrm{LLM}}^2 L)$ for a full LLM forward pass.  The LLM reads
$r_t^{(K)}$ instead of $r_t^{(0)}$, enabling cheap amortised reasoning.

% ---------------------------------------------------------------------------
\section{Experimental Setup}
\label{sec:setup}
% ---------------------------------------------------------------------------

\subsection{Models}

\begin{table}[h]
\centering
\caption{Models and configurations.}
\label{tab:models}
\begin{tabular}{@{}llll@{}}
\toprule
Label & Base Model & Memory Component & Training \\
\midrule
Control 1 & Qwen3.5-0.8B (frozen) & RoPE / YaRN & Zero \\
Control 2 & Mamba-2 (1.3B) & Native SSM & Zero \\
Control 3 & Qwen3.5-0.8B + LoRA & Infini-attention & Low \\
Track A-ro & Qwen3.5-0.8B + LoRA & ESN sidecar (read-only) & Low \\
Track A-rw & Qwen3.5-0.8B + LoRA & ESN sidecar (read/write) & Low \\
Track A-sweep & Qwen3.5-0.8B + LoRA & Best HP from T16 & Low \\
Arch Control & LLaMA-3.2-1B + LoRA & ESN sidecar & Low \\
\bottomrule
\end{tabular}
\end{table}

All Track A/Arch Control models use LoRA rank 16 ($\alpha=32$) on
$\{q, k, v, o, \text{gate}, \text{up}, \text{down}\}$ projections plus
the read/write interface parameters.  CPU-only training uses float32;
GPU training uses bfloat16 with gradient checkpointing.

\subsection{Benchmarks}
\label{sec:benchmarks}

\paragraph{Passkey Retrieval.}
Synthetic long-context task: a ``passkey'' (5-digit number) is buried in
random text; the model must retrieve it.  Evaluated at context lengths
$\{1k, 2k, 4k, 8k, 16k, 32k\}$ tokens.  Metric: exact-match accuracy.

\paragraph{Variable Tracking.}
A sequence of variable assignments and updates is presented; the model must
report the final value of a queried variable.  Evaluates algorithmic
working memory.

\paragraph{Associative Recall.}
Key-value pairs are presented in a list, followed by a query for one key.
The model must return the correct value.  Tests content-addressable retrieval.

\paragraph{Compositional Generalization.}
Held-out combinations of operator types and argument lengths, testing
whether the model generalises beyond its training distribution.

\paragraph{Chaos Diagnostics.}
Small-scale chaotic system forecasting (Lorenz attractor) as a representation
diagnostic.  Verifies that reservoir states exhibit stable bounded dynamics.

\paragraph{Perplexity.}
FineWeb held-out text perplexity, measuring language modelling degradation.

\subsection{Training Details}

Optimizer: AdamW with learning rate $2 \times 10^{-4}$ (LoRA) and $10^{-3}$
(interface), weight decay $0.01$, gradient clip $1.0$.  Schedule: cosine
decay with 100 warmup steps over 5000 total steps.  Batch size 4 with
gradient accumulation $\times 4$ (effective batch 16).  Data: 60\% FineWeb,
30\% synthetic task data, 10\% chaotic systems sequences.

% ---------------------------------------------------------------------------
\section{Results}
\label{sec:results}
% ---------------------------------------------------------------------------

\subsection{Gate A Status}

\begin{table}[h]
\centering
\caption{Gate A criteria and current status (all GPU runs pending).}
\label{tab:gate_a}
\begin{tabular}{@{}lllll@{}}
\toprule
\# & Criterion & Threshold & Best value & Result \\
\midrule
1 & Passkey retrieval gain & $\geq 10\%$ & Pending & \textbf{PENDING} \\
2 & Algorithmic memory gain & $\geq 15\%$ & Pending & \textbf{PENDING} \\
3 & Compositional gen.\ gain & $\geq 10\%$ & Pending & \textbf{PENDING} \\
4 & Inference latency overhead & $\leq 20\%$ & Pending & \textbf{PENDING} \\
5 & Perplexity degradation & $< 2\%$ & Pending & \textbf{PENDING} \\
\bottomrule
\end{tabular}
\end{table}

All Gate~A results require GPU training to complete.
Table~\ref{tab:gate_a} will be updated with empirical results.
Per our pre-registered negative-result commitment, if any gate criterion is
not met, the paper will include a detailed failure analysis in
Section~\ref{sec:discussion}.

\subsection{Baseline Comparison}

\begin{table}[h]
\centering
\caption{Full model comparison (values to be filled from GPU runs).}
\label{tab:full_results}
\begin{tabular}{@{}lrrrrr@{}}
\toprule
Model & Passkey $\uparrow$ & Algo Mem $\uparrow$ & Comp Gen $\uparrow$
      & Latency $\downarrow$ & PPL $\downarrow$ \\
\midrule
Control 1 (Qwen+YaRN)  & --- & --- & --- & --- & --- \\
Control 2 (Mamba-2)    & --- & --- & --- & --- & --- \\
Control 3 (Infini-attn)& --- & --- & --- & --- & --- \\
Track A read-only      & --- & --- & --- & --- & --- \\
Track A read/write     & --- & --- & --- & --- & --- \\
Track A best sweep     & --- & --- & --- & --- & --- \\
Arch Control (LLaMA+RC)& --- & --- & --- & --- & --- \\
\bottomrule
\end{tabular}
\end{table}

\subsection{HP Sweep Results (Track A)}

Figure~\ref{fig:pareto} shows the quality-vs-efficiency Pareto frontier from
the reservoir hyperparameter sweep (T16).  Table~\ref{tab:sweep_summary}
summarises the effect of each hyperparameter.

\begin{figure}[h]
\centering
\includegraphics[width=0.7\textwidth]{figures/pareto_frontier.pdf}
\caption{Quality vs.\ inference latency overhead Pareto frontier for Track~A
  reservoir hyperparameter sweep.  Each point is a configuration from
  Table~\ref{tab:sweep_configs}.  Dashed lines mark the Gate~A thresholds.
  \textit{(Requires GPU runs to populate.)}}
\label{fig:pareto}
\end{figure}

\subsection{Ablation Analysis}
\label{sec:ablations}

Table~\ref{tab:ablations} summarises ablation results (from T29 scripts).
Key ablations: (i) read-only vs.\ read/write, (ii) single vs.\ multi-reservoir,
(iii) randomised/stateless control (verifies recurrent dynamics, not just
additional parameters, drive gains), (iv) spectral radius regimes.

\begin{table}[h]
\centering
\caption{Ablation summary. $\Delta$ values relative to full read/write model.}
\label{tab:ablations}
\begin{tabular}{@{}lrrrr@{}}
\toprule
Ablation & $\Delta$ Passkey & $\Delta$ Algo Mem & $\Delta$ PPL & Notes \\
\midrule
Read-only (no write) & --- & --- & --- & Interface ablation \\
Randomised state (no ESN dynamics) & --- & --- & --- & Control \\
Single reservoir (fast only) & --- & --- & --- & Multi-res ablation \\
Single reservoir (slow only) & --- & --- & --- & Multi-res ablation \\
Spectral radius $\rho=0.5$ & --- & --- & --- & Stability regime \\
Spectral radius $\rho=1.1$ & --- & --- & --- & Near-chaotic \\
\bottomrule
\end{tabular}
\end{table}

% ---------------------------------------------------------------------------
\section{Discussion}
\label{sec:discussion}
% ---------------------------------------------------------------------------

\paragraph{DeltaNet synergy.}
Qwen3.5's Gated DeltaNet blocks use a delta-rule update $\Delta h_t = \beta
(\hat{v}_t - h_{t-1} \cdot \hat{k}_t^\top \hat{k}_t)$ analogous to the ESN
leaky integration (Eq.~\ref{eq:esn}).  We hypothesise that this shared
inductive bias reduces the distributional mismatch between LLM hidden states
and ESN reservoir states, making the read/write interface easier to train.
The Arch Control (LLaMA-3.2-1B + ESN, pure softmax attention) provides the
counterfactual.  A larger $\Delta$ improvement for Qwen3.5 vs.\ LLaMA would
support the DeltaNet synergy hypothesis.

\paragraph{Latent reasoning cost analysis.}
A single ESN sub-step (Eq.~\ref{eq:substep}) costs
$O(\rho_s D_r^2)$ FLOPs where $\rho_s \approx 0.01$ is reservoir sparsity.
For $D_r = 10{,}000$: $\sim 10^6$ FLOPs/step.  A Qwen3.5-0.8B forward pass
costs $\sim 1.6 \times 10^9$ FLOPs.  Thus $K=16$ sub-steps cost $\sim 1\%$
of a single LLM forward pass---making reservoir latent reasoning
approximately $100\times$ cheaper per reasoning step than Coconut-style
hidden-state recirculation.

\paragraph{What worked.}
[To be filled after GPU runs.]

\paragraph{What did not work / failure analysis.}
[To be filled after GPU runs.  Per our pre-registered commitment, if Gate~A
criteria are not met, we document: (i) which criterion failed and by how much,
(ii) diagnostic experiments to understand the failure, (iii) recommendations
for future work.]

\paragraph{Reservoir memory saturation.}
At long context lengths ($L \gg 1/\alpha$), the reservoir state integrates all
past inputs and may saturate.  The dual-timescale design (Track~B) partially
addresses this: the fast reservoir refreshes quickly while the slow reservoir
maintains long-range history.  For very long contexts ($L > 10^5$), periodic
reservoir resets or explicit memory management may be needed.

% ---------------------------------------------------------------------------
\section{Conclusion}
\label{sec:conclusion}
% ---------------------------------------------------------------------------

We have presented the Latent Reservoir Scratchpad (\LRS), an approach to
augmenting LLMs with a fixed-size continuous-state working memory derived from
Echo State Networks.  Key advantages are $O(1)$ memory footprint,
interpretable dynamics (the reservoir weights are known and fixed), and
extremely cheap latent reasoning sub-steps ($\sim 100\times$ cheaper than
LLM-internal recirculation).

The three-track evaluation program (bolt-on sidecar, inserted layers,
from-scratch hybrid) with quantitative success gates and a pre-registered
negative-result commitment provides a rigorous empirical foundation for
assessing the viability of reservoir-augmented LLMs.  GPU evaluation results
will be incorporated as they become available.

Code, configs, and reproducibility instructions are available at
\url{https://github.com/[repository-url]}.

% ---------------------------------------------------------------------------
\bibliography{references}
\bibliographystyle{plainnat}
% ---------------------------------------------------------------------------

\appendix

% ---------------------------------------------------------------------------
\section{Hyperparameter Sweep Configurations}
\label{app:sweep}
% ---------------------------------------------------------------------------

\begin{table}[h]
\centering
\caption{Reservoir HP sweep grid (T16). Best combo selected by Pareto
  criterion: max passkey accuracy at $\leq 10\%$ latency overhead.}
\label{tab:sweep_configs}
\begin{tabular}{@{}lllll@{}}
\toprule
Config & $D_r$ & $\alpha$ & $\rho$ & Topology \\
\midrule
size-500 & 500 & 0.3 & 0.9 & erdos\_renyi \\
size-2000 & 2000 & 0.3 & 0.9 & erdos\_renyi \\
size-10000 & 10000 & 0.3 & 0.9 & erdos\_renyi \\
size-50000 & 50000 & 0.3 & 0.9 & erdos\_renyi \\
leak-0.1 & 1000 & 0.1 & 0.9 & erdos\_renyi \\
leak-0.3 & 1000 & 0.3 & 0.9 & erdos\_renyi \\
leak-0.7 & 1000 & 0.7 & 0.9 & erdos\_renyi \\
leak-1.0 & 1000 & 1.0 & 0.9 & erdos\_renyi \\
sr-0.5 & 1000 & 0.3 & 0.5 & erdos\_renyi \\
sr-0.9 & 1000 & 0.3 & 0.9 & erdos\_renyi \\
sr-0.99 & 1000 & 0.3 & 0.99 & erdos\_renyi \\
sr-1.1 & 1000 & 0.3 & 1.1 & erdos\_renyi \\
small-world & 1000 & 0.3 & 0.9 & small\_world \\
best-combo & 2000 & 0.3 & 0.9 & small\_world \\
\bottomrule
\end{tabular}
\end{table}

% ---------------------------------------------------------------------------
\section{Architecture Diagrams}
\label{app:arch}
% ---------------------------------------------------------------------------

Figures~\ref{fig:sidecar_arch}--\ref{fig:rw_transformer} illustrate the
three integration depths.  These figures are generated by the scripts in
\texttt{docs/paper/figures/}.

\begin{figure}[h]
\centering
\includegraphics[width=0.85\textwidth]{figures/architecture_sidecar.pdf}
\caption{Track~A bolt-on reservoir sidecar.  Frozen LLM (left) passes hidden
  states to the WriteHead; ESN state is projected back via ReadProjection and
  injected as K/V in the CrossAttentionSidecar.}
\label{fig:sidecar_arch}
\end{figure}

\begin{figure}[h]
\centering
\includegraphics[width=0.85\textwidth]{figures/architecture_ril.pdf}
\caption{Track~B Reservoir Interaction Layer (RIL).  Dual-timescale
  multi-reservoir replaces selected DeltaNet blocks; FiLM modulation
  injects reservoir state into residual stream.}
\label{fig:ril_arch}
\end{figure}

\begin{figure}[h]
\centering
\includegraphics[width=0.85\textwidth]{figures/architecture_rw_transformer.pdf}
\caption{Track~C RW-Transformer decoder block.  Three parallel branches
  (attention, MLP, dual-reservoir workspace) fused by a learned gating layer.}
\label{fig:rw_transformer}
\end{figure}

\end{document}
